\newpage
\section{مدل های \lr{Energy-Based}}
\subsection{سوال اول}
روش \lr{Contrastive Divergence} راهی برای تخمین گرادیان لگاریتم تابع درستنمایی در حین آموزش مدل های \lr{Energy-based} هست.
در این الگوریتم زنجیره \lr{MCMC} را از یک نقطه $x$  شروع میکنیم و برای تعدادی گام مشخص ادامه میدهیم. در نهایت میخواهیم تابع هدف زیر را بهینه کنیم :
\begin{equation*}
    CD_T = D_{\mathbb{KL}} (p_0 ||p_\infty) - D_{\mathbb{KL}}(p_T ||p_\infty)
\end{equation*}
که در اینجا $p_T$ توزیع تخمین زده شده برای داده $x$ با استفاده از \lr{MCMC} بعد از $T$  قدم هست و $p_0$ توزیع داده امون هست.\\
دلیل زمانبر بودن با اینکه حتی برای رسیدن به \lr{equilibrium} صبر نمیکنیم میتواند به دلایل زیر باشد :
\begin{enumerate}
    \item در داخل الگوریتم \lr{MCMC} یک حلقه سمپلینگ گیبز یا \lr{Langevin Dynamics} نیز قرار میدهیم تا بتوانیم از سمپل های منفی هم بهره ببریم در آموزش. همین باعث کند شدن الگوریتم میشود.
    \item در ابعاد بالا حتی برای چند قدم ، اجرا کردن این الگوریتم میتواند از نظر محاسباتی زمانبر باشد.
\end{enumerate}
\subsection{سوال دوم}

یک مزیت نسبت به \lr{Metropolis-Hastings} این هست که سمپل های $i.i.d$ تولید میکند در حالی که در $MH$ یک زنجیره مارکف داریم که هر سمپل وابسته به سمپل قبل از خود هست.

یک نقطه ضعف نسبت به الگوریتم $MH$ این هست که \lr{rejection sampling} از نفرین ابعاد رنج میبرد. در واقع برای ابعاد بسیار بالا پیدا کردن یک توزیع $q(x)$ و $M$ مناسب که در رابطه $p(x)\leq Mq(x)$صدق کند به شدت سخت هست چون به ازای انتخاب های بزرگ $M$ احتمال قبولی متناسب با $1/M$ میشود و تقریبا تمام سمپل های پیشنهادی ریجکت میشوند.

\subsection{سوال سوم}
با توجه به توزیع داده شده 
\begin{equation*}
    p(x) = 0.6 \mathcal{N}(50,1) + 0.4\mathcal{N}(0.001,\sigma^2) 
\end{equation*}
\begin{equation*}
    Q(x'|x) = \mathcal{N}(x,0.5) 
\end{equation*}
دلایل کند بودن الگوریتم $MH$ برای سمپل کردن از توزیع بالا :
\begin{enumerate}
    \item \lr{Mode Isolation} : مد ها در $x\approx 0$ و $x\approx 50$ قرار گرفته اند. توزیع پیشنهادی استپ سایز بسیار کوچکی دارد . احتمال اینکه الگوریتم $MH$ یک جامپ از روی ۵۰ واحد فاصله پیشنهاد دهد تا به مد دوم برسیم تقریبا صفر هست.
    \item \lr{Poor Mixing} : این زنیجره احتمالا در هر مدی که از آن شروع کند گیر میکند و همین باعث \lr{Poor Mixing} میشود.
\end{enumerate}
راه حل های پیشنهادی عبارتند از 
\begin{enumerate}
    \item \lr{Parallel Tempering} : اجرای چندین زنجیره متفاوت با \lr{temperature} های مختلف تا فضای احتمالی را صاف کنیم.
    \item از یک مدل پیشنهادی استفاده کنیم که به طور معمول از استپ سایز های بزرگ استفاده میکند تا مانع گیر کردن در یک مد شویم.
\end{enumerate}